\clearpage
\section{Anomaly Detection in Signals}
\subsection{Classical Statistical Methods For Anomaly Detection In Time Series Signals}
We'll break this question down into the different methods mentioned and explain how each one can be used for anomaly detection. All of the methods assume there is an expected behavior and the anomalies deviate from the expected behavior.
\subsubsection{Thresholding}
If we assume the signal has some median value and a certain distribution around it, we can set an upper bound on the deviation from the median that's considered acceptable. Values exceeding this threshold will be flagged as anomalies. This is very simple to implement but its downsides are obvious: the assumption that the signal is well behaved around some median and doesn't spike too often.
\subsubsection{Spectral Analysis}
If we know our signal is expected to be comprised of a specific set or range of frequencies, we can perform a fourier transform on it and in the frequency domain we will have an acceptable range or set of frequencies, with any frequency outside this range or set being considered anomalous. This method is useful for signals which either are expected to be periodic, carried by specific frequencies, or are specifically characterized by their frequency. The downsides to this method is the method by which the signal is analyzed itself may introduce artifacts or lose information, non-repeating anomalies may be hard to detect, very slow changes in the signal may be hard to detect, and of course if the signal isn't strongly characterized by its frequency this method won't be effective.
\subsubsection{Kalman Filtering Residuals}
Kalman filtering gives an optimal linear estimator for a signal. By assuming a model for the signal, we can use a Kalman filter to predict the next value of the signal. By monitoring the difference between the predicted value and the measured value and setting a threshold on this difference, we can declare measurements which exceed this threshold (which may be dynamic or static) as anomalies. This method is useful when we have a good understanding of the system generating this data and can model it well, and the system in question behaves in a linear manner. The downsides to Kalm Filtering are that it assumes the system is linear in its evolution, and that we have a good model of the system.
\subsection{Encoder-Decoder For Anomaly Detection}
If we treat an incoming signal as a \(d_1\) dimensional vector but we know the underlying signal is of a lower dimensionality, we can train a neutral network to compress this vector into a \(d_2 < d_1\) dimensional vector, then expand it back to a \(d_1\) dimensional vector. We can then train this network on known good data, and the NN will learn to reconstruct the signal. 